\documentclass[10pt,a4paper]{amsart}
\usepackage[margin=19mm]{geometry}
\usepackage{amsmath,amssymb,mathtools}
\usepackage[scr=rsfs,cal=euler]{mathalfa}
\usepackage{xfrac}
\usepackage[unicode,hidelinks,bookmarks=false]{hyperref}
\newcommand{\ie}{i.e.\ }
\newcommand{\cf}{cf.\ }
\newcommand{\resp}{respectively\ }
\newcommand{\Subsection}{\subsection}
\providecommand{\nobreakdash}{\nobreak\hbox{-}}
\setlength{\emergencystretch}{2em}
\multlinegap0pt
\usepackage{amssymb}
\usepackage[all]{xy}
\newtheorem{thm}{Theorem}
\title{A generalisation of the Burau representation and groups $G_{n}^{3}$ for classical braids}
\author{Vassily Olegovich Manturov, Igor Mikhailovich Nikonov}
\date{Source: arXiv:2602.07890v1 (CC BY 4.0)}
\begin{document}
\maketitle

\noindent\emph{Source and licence.} A generalisation of the Burau representation and groups $G_{n}^{3}$ for classical braids. Version arXiv:2602.07890v1 (CC BY 4.0), distributed by arXiv under the Creative Commons Attribution 4.0 International licence, http://creativecommons.org/licenses/by/4.0/, as recorded in the arXiv licence metadata for this exact version and shown on the abstract page https://arxiv.org/abs/2602.07890v1. Full licence notice: \texttt{licenses/arxiv-2602.07890-CC-BY-4.0.txt}.\par\smallskip

\noindent\textit{Editorial scope.} Counted unit: the article's thm thm (source0.tex lines 548--550). The article's complete original proof of it is delivered with it (source0.tex lines 552--567, one \texttt{proof} environment of 16 lines). No mathematics is written, repaired or re-derived: the statement and the proof are the article's own lines, in the article's own notation, and no retained line is substituted or re-flowed.  No further source-local context is needed: every cross-reference of every retained line resolves inside the two delivered spans.  Nothing was omitted from any retained span, so the package carries no omission record.  No retained line contains a citation, so the wrapper carries no bibliography and no bibliography entry.  No retained line contains a figure environment, an image-inclusion command or drawing code, so no image is delivered and the package declares no asset dependency.  Editorial formatting only, in addition: an AMS \texttt{amsart} preamble that restates the article's own declarations and macros used by the retained lines, namely the environments \texttt{thm} on separate counters, together with the standard packages those lines need.  The retained environments print in this document as Theorem 1 in order of appearance, whereas the article prints the same items with the numbering of its own full document.  The complete native source of the article as distributed in the arXiv e-print for this exact version is bundled unchanged as \texttt{sources/194219/source0.tex}.
\begin{thm}
    The map $w$ is invariant under first and second Markov moves.
\end{thm}

\begin{proof}
    1. Consider a first Markov move $\beta\mapsto\beta'=\beta_1\beta\beta_1^{-1}$. Denote $y=\psi\circ\phi(\tilde\beta)$, $y_1=\psi\circ\phi(\tilde\beta_1)\in\Sigma_n\ltimes ZH$. Then
\[
w(\beta')-w(\beta)=p^{c\circ\sigma_1^{-1}}(y_1yy_1^{-1})-p^c(y)=0\in\bar A_C,    
\]
where $c=\tilde c_\beta$ and $\sigma_1=\rho(\tilde\beta_1)$.

2. Let $\beta\mapsto \beta'=\sigma_n\beta\in B_{n+1}$ be a second Markov move. Then $c_{\beta'}(i)=c_\beta(i)$ for $i\le n$ and $c_{\beta'}(n+1)=c_\beta(n)$.

Let $v=\phi(\tilde\beta)\in G^3_{2n}$ and $v'=\phi(\tilde\beta')\in G^3_{2n+2}$. Then $v'=v'_1v'_2$ where $v'_1$ corresponds to $\sigma_n$ and $v'_2$ corresponds to $\beta$.

Since the $v'_1$ consists of generators $a_{n,n+1,j}$ and the impact of a generator $a_{ijk}$ to the image of the map $\psi$ is nontrivial only if the components $c(i)$, $c(j)$ and $c(k)$ are all different, we can ignore the term $v'_1$.

The word $v'_2$ differs from the word $v$ by the shift by $1$ of the indices which are greater than $n$, and by 2-letter subwords $a_{ij,n+1}a_{ij,n+2}$ which may appear in some places of $v'_2$. Since the strands $n+1$ and $n+2$. belong to one component, the impact of these generators in $Z_{c_1c_2c_3}$ by the map $p^c$ annihilates. Hence, $p^c(v)=p^c(v'_2)=p^c(v')$.
Thus, $w(\beta)=w(\beta')$.
\end{proof}

\end{document}
